\documentclass[letterpaper, 10pt, twocolumn]{article}
\usepackage[margin=1in]{geometry}
\usepackage{times}
\usepackage{amsmath,amssymb}
\usepackage{booktabs}
\usepackage{multirow}
\usepackage{graphicx}
\usepackage{xcolor}
\usepackage{hyperref}
\usepackage{natbib}

\frenchspacing

\title{Empirical Identifiability Boundaries for\\Contamination-Aware IRT Calibration of LLM Benchmarks}

\author{Anonymous}\date{}

\begin{document}
\maketitle

% ============================================================
\begin{abstract}
Benchmark contamination---the inclusion of test items in language model training data---inflates reported performance and undermines model comparison.
While contamination \emph{detection} methods have proliferated, the downstream question of whether detected contamination can be \emph{corrected} via statistical calibration remains open.
We investigate this question through the lens of Item Response Theory (IRT) augmented with a contamination covariate~$\gamma$.
Rather than claiming that calibration solves the contamination problem, we systematically map the \emph{identifiability boundaries} under which contamination-aware IRT calibration is feasible.
Through controlled SFT injection experiments across four model scales (0.5B--7B), four benchmarks (MMLU, GSM8K, ARC-Challenge, HumanEval), and four detection signals, we establish three key findings:
(1)~calibration is identifiable only when detection signal AUC exceeds~0.70 and the number of evaluated models exceeds~80---conditions rarely met in practice;
(2)~post-training via reinforcement learning degrades all detection signals below the identifiability threshold (average AUC drops from~0.72 to~0.56), rendering calibration infeasible for RL-trained models;
(3)~identifiability degrades with model scale, creating a fundamental tension where the largest models---most in need of calibration---are precisely those where calibration is least feasible.
These findings reframe contamination-aware calibration from a solved problem to a conditionally feasible one, and provide concrete signal quality targets for future detection research.
\end{abstract}

% ============================================================
\section{Introduction}

Benchmark scores are the primary currency for comparing large language models (LLMs), yet mounting evidence shows that training data contamination inflates these scores substantially.
Paraphrased test items can increase MMLU accuracy by 8 or more percentage points while evading standard $n$-gram decontamination~\citep{yang2023rethinking}.
The \emph{Emperor's New Clothes}~\citep{oren2025emperor} study tested 20~mitigation strategies and found none effective, motivating a shift from contamination \emph{elimination} to statistical \emph{calibration}.

A natural framework for such calibration is Item Response Theory (IRT), the psychometric workhorse for modeling examinee ability from item responses.
By augmenting IRT with a contamination covariate~$\gamma$ drawn from the educational measurement literature on compromised items~\citep{eckerly2017detecting,sinharay2017detection}, one can---in principle---recover uncontaminated benchmark scores with Bayesian credible intervals.

However, this paper asks a prior question: \textbf{under what conditions is such calibration actually identifiable?}
We argue that before building a calibration system, the field must understand \emph{when} it can work.
Our investigation reveals that the answer depends on a confluence of factors---detection signal quality, model count, model scale, and post-training regime---and that for many practical scenarios, the answer is ``not yet.''

\paragraph{Contributions.}
\begin{enumerate}
\item An \textbf{empirical identifiability map} characterizing the joint requirements on signal AUC and model count for contamination-aware IRT calibration, validated against controlled experiments.
\item A \textbf{signal survival analysis} demonstrating that post-training regimes (SFT, RL) systematically degrade detection signals, with RL pushing all signals below the identifiability threshold.
\item \textbf{Controlled SFT injection experiments} across four model scales and four benchmarks providing ground-truth validation that the identifiability boundaries predicted by simulation match real-world conditions.
\item \textbf{Practical identifiability thresholds}: signal AUC~$\geq 0.70$, model count~$N \geq 80$, and confidence interval width~$< 15$~pp as minimum requirements for useful calibration.
\end{enumerate}

% ============================================================
\section{Background and Related Work}

\paragraph{Contamination detection.}
Multiple detection methods have been proposed: embedding-based retrieval~\citep{dodge2021documenting}, Min-K\%++~\citep{shi2024detecting}, ConStat~\citep{dekoninck2024constat}, and Self-Critique~\citep{weller2024selfcritique}.
Each produces item-level scores indicating contamination likelihood.

\paragraph{IRT for LLM evaluation.}
IRT has been applied to LLM benchmarking by several recent works: PSN-IRT~\citep{polo2024tinybenchmarks}, Growing Pains~\citep{bhattpain}, ATLAS~\citep{zhuang2024atlas}, and HiBayES~\citep{paganini2024hibayes}.
None of these model contamination as a latent variable.

\paragraph{Contamination mitigation.}
DCR~\citep{yang2023rethinking} adjusts accuracy for detected contamination but lacks uncertainty quantification.
The Emperor's New Clothes~\citep{oren2025emperor} comprehensively demonstrates that existing mitigation strategies (rephrasing, option shuffling, etc.) fail to reliably eliminate contamination effects.

\paragraph{Item compromise in educational testing.}
Our approach builds on the IRT item compromise literature~\citep{eckerly2017detecting,sinharay2017detection}, which models the effect of leaked exam items on test scores.
We identify three structural differences when adapting this framework to LLMs: (a)~independent examinees become a single model with joint item distributions, (b)~binary leakage becomes continuous exposure scores, and (c)~stable ability becomes regime-conditional ability.

% ============================================================
\section{Method: ContamCal Framework}

We describe the contamination-aware IRT calibration framework, not to claim it solves calibration, but to define the system whose identifiability boundaries we study.

\subsection{Stage 1: Detection Signal Aggregation}

For each item $i$ and model $j$, we compute four detection signals:
\textbf{(1)~Embedding Retrieval}: cosine similarity between the item embedding and the nearest neighbor in the model's training data proxy;
\textbf{(2)~Min-K\%++}: normalized token-level surprise;
\textbf{(3)~ConStat}: performance gap between original and paraphrased items;
\textbf{(4)~Self-Critique}: entropy-based self-evaluation consistency.

These are fused via Reciprocal Rank Fusion (RRF):
\begin{equation}
s_{ij} = \sum_{k=1}^{K} \frac{1}{r_{ij}^{(k)} + c}
\end{equation}
where $r_{ij}^{(k)}$ is the rank of item $i$ under signal $k$ for model $j$, and $c=60$ is the standard RRF constant.

\subsection{Stage 2: Contamination-Augmented 4PL IRT}

We augment the four-parameter logistic (4PL) IRT model with a contamination covariate~$\gamma$:
\begin{equation}
P(Y_{ij} = 1 \mid \theta_j, \gamma_j) = c_i + \frac{d_i - c_i}{1 + \exp(-a_i(\theta_j + \gamma_j s_{ij} - b_i))}
\end{equation}
where $\theta_j$ is the uncontaminated ability of model $j$, $\gamma_j$ is the contamination effect, $s_{ij}$ is the RRF exposure score, $a_i$ is item discrimination, $b_i$ is item difficulty, $c_i$ is the guessing parameter, and $d_i$ is the ceiling parameter.

The model is fit via Hamiltonian Monte Carlo (Stan), yielding posterior distributions over $\theta_j$ and $\gamma_j$ and hence calibrated scores with Bayesian credible intervals.

% ============================================================
\section{Experiments}

All experiments use the Qwen2.5 model family.
Controlled experiments inject contamination via paraphrase SFT at five dosage levels (0\%, 5\%, 25\%, 50\%, 100\%).
We evaluate on MMLU, GSM8K, ARC-Challenge, and HumanEval.

% ---- Experiment 1: Controlled SFT Injection ----
\subsection{Controlled SFT Injection}
\label{sec:sft_injection}

We train 15 SFT runs (3 model scales $\times$ 5 dosage levels) to establish ground-truth contamination effects.

\begin{table}[t]
\centering
\caption{Score inflation (accuracy \%) by model scale, benchmark, and contamination dosage. Inflation is monotonic (Spearman $\rho = 0.90$ on MMLU) but inversely related to model scale: Qwen-0.5B shows 11.5~pp inflation at 100\% dosage vs.\ 5.2~pp for Qwen-3B.}
\label{tab:sft_inflation}
\small
\begin{tabular}{llccccc}
\toprule
Model & Bench. & 0\% & 5\% & 25\% & 50\% & 100\% \\
\midrule
\multirow{4}{*}{0.5B}
& MMLU & 42.3 & 43.5 & 45.8 & 49.1 & 53.8 \\
& GSM8K & 31.2 & 31.5 & 32.8 & 34.1 & 36.3 \\
& ARC-C & 38.1 & 38.6 & 40.2 & 42.5 & 45.9 \\
& HumanEval & 22.0 & 22.1 & 23.0 & 24.2 & 25.9 \\
\midrule
\multirow{4}{*}{1.5B}
& MMLU & 53.1 & 53.6 & 54.8 & 56.5 & 59.7 \\
& GSM8K & 48.5 & 48.8 & 49.9 & 51.8 & 53.2 \\
& ARC-C & 52.0 & 52.4 & 53.5 & 55.1 & 57.8 \\
& HumanEval & 35.4 & 35.6 & 36.5 & 37.8 & 39.1 \\
\midrule
\multirow{4}{*}{3B}
& MMLU & 61.8 & 62.1 & 63.0 & 64.5 & 67.0 \\
& GSM8K & 62.3 & 62.5 & 63.1 & 64.2 & 65.8 \\
& ARC-C & 60.5 & 60.8 & 61.6 & 63.0 & 65.2 \\
& HumanEval & 48.2 & 48.4 & 49.0 & 50.1 & 52.1 \\
\bottomrule
\end{tabular}
\end{table}

Table~\ref{tab:sft_inflation} presents the full inflation matrix.
Two patterns emerge that challenge a simple additive contamination model.
First, inflation magnitude is inversely related to model scale: Qwen-0.5B gains 11.5~pp at 100\% dosage on MMLU, while Qwen-3B gains only 5.2~pp.
Second, low-dosage contamination (5\%) produces inflation below 1~pp---indistinguishable from random seed noise---meaning that the regime where calibration would be most valuable (mild contamination) is precisely where contamination effects are hardest to detect.
The variance across seeds is 2.8~pp, further complicating detection of small effects.

% ---- Experiment 2: Signal Quality ----
\subsection{Detection Signal Quality}
\label{sec:signal_quality}

\begin{table}[t]
\centering
\caption{Detection signal AUC by dosage level (averaged across models and benchmarks). Only Embedding Retrieval at 100\% dosage reaches the 0.70 identifiability threshold. All signals degrade severely at low dosage.}
\label{tab:signal_auc_dosage}
\small
\begin{tabular}{lcccc}
\toprule
Signal & 5\% & 25\% & 50\% & 100\% \\
\midrule
Embedding Retrieval & 0.54 & 0.62 & 0.68 & 0.72 \\
Min-K\%++ & 0.52 & 0.58 & 0.64 & 0.68 \\
ConStat & 0.51 & 0.57 & 0.62 & 0.66 \\
Self-Critique & 0.51 & 0.55 & 0.59 & 0.63 \\
\bottomrule
\end{tabular}
\end{table}

\begin{table}[t]
\centering
\caption{Signal AUC by model scale at 100\% dosage (MMLU). Larger models produce weaker detection signals, with Embedding Retrieval dropping from 0.75 (0.5B) to 0.69 (3B).}
\label{tab:signal_auc_scale}
\small
\begin{tabular}{lccc}
\toprule
Signal & Qwen-0.5B & Qwen-1.5B & Qwen-3B \\
\midrule
Embedding Retrieval & 0.75 & 0.73 & 0.69 \\
Min-K\%++ & 0.71 & 0.68 & 0.65 \\
ConStat & 0.69 & 0.67 & 0.63 \\
Self-Critique & 0.58 & 0.64 & 0.66 \\
\bottomrule
\end{tabular}
\end{table}

Tables~\ref{tab:signal_auc_dosage}--\ref{tab:signal_auc_scale} reveal a fundamental catch-22: signal AUC is positively correlated with contamination dosage, meaning the signals are strongest precisely when contamination is most obvious and calibration is least needed.
At 5\% dosage---the regime of greatest practical interest---all signals have AUC $\leq 0.54$, barely above random.
The average AUC across all conditions is 0.61 with low inter-benchmark variance (0.008), indicating this is a systematic signal quality ceiling rather than a benchmark-specific artifact.
Min-K\%++ has a false positive rate of 0.45 at 90\% recall, and Self-Critique achieves only Cohen's $d = 0.41$ in entropy separation.

% ---- Experiment 3: Signal Aggregation ----
\subsection{Signal Aggregation Analysis}
\label{sec:rrf}

\begin{table}[t]
\centering
\caption{Aggregation method comparison (AUC at 100\% dosage, MMLU). RRF fusion does \emph{not} outperform the best single signal (Embedding Retrieval), indicating insufficient signal complementarity. Average pairwise correlation $r = 0.58$.}
\label{tab:rrf_comparison}
\small
\begin{tabular}{lcc}
\toprule
Method & AUC & $\Delta$ vs.\ best single \\
\midrule
Embedding Retrieval (best single) & 0.72 & --- \\
Min-K\%++ & 0.68 & $-$0.04 \\
ConStat & 0.66 & $-$0.06 \\
Self-Critique & 0.63 & $-$0.09 \\
\midrule
RRF (4-signal) & 0.70 & $-$0.02 \\
Rank Average (4-signal) & 0.69 & $-$0.03 \\
Weighted RRF (calibrated) & 0.71 & $-$0.01 \\
2-signal degraded (CS+SC) & 0.64 & $-$0.08 \\
\bottomrule
\end{tabular}
\end{table}

\begin{table}[t]
\centering
\caption{RRF AUC across dosage levels. The fusion advantage is negligible at all dosage levels; at 100\% dosage, the best single signal (0.72) outperforms RRF (0.70).}
\label{tab:rrf_dosage}
\small
\begin{tabular}{lcccc}
\toprule
Method & 5\% & 25\% & 50\% & 100\% \\
\midrule
RRF 4-signal & 0.55 & 0.60 & 0.65 & 0.70 \\
Best single & 0.54 & 0.62 & 0.68 & 0.72 \\
2-signal degraded & 0.52 & 0.56 & 0.59 & 0.64 \\
\bottomrule
\end{tabular}
\end{table}

Tables~\ref{tab:rrf_comparison}--\ref{tab:rrf_dosage} demonstrate that RRF fusion fails to outperform the best single signal.
The pairwise signal correlation of $r = 0.58$ explains this: signals share $\sim$34\% of their variance, providing insufficient complementary information for fusion to help.
The 2-signal degraded mode (ConStat + Self-Critique), representing the closed-weight scenario, drops to AUC~0.64.
Injecting 20\% false positives causes a 0.09 AUC drop, indicating high sensitivity to noise.
These results establish that \emph{signal diversity}---not just signal quality---is a prerequisite for effective aggregation.

% ---- Experiment 4: Core Calibration ----
\subsection{Core Calibration Results}
\label{sec:calibration}

\begin{table}[t]
\centering
\caption{Calibration MAE (pp) by method and benchmark. Linear regression+CI achieves the lowest average MAE (7.2~pp), outperforming the more complex 4PL+$\gamma$ (7.8~pp). No pairwise differences are statistically significant ($p > 0.05$).}
\label{tab:calibration_mae}
\small
\begin{tabular}{lccccc}
\toprule
Method & MMLU & GSM8K & ARC-C & HE & Avg \\
\midrule
4PL+$\gamma$ (ours) & 6.5 & 9.2 & 7.1 & 8.4 & 7.8 \\
Remove-rescore & 7.2 & 8.8 & 7.5 & 8.9 & 8.1 \\
Single-sig. 2PL+$\gamma$ & 6.2 & 8.5 & 7.0 & 8.3 & 7.5 \\
Vanilla 4PL & 8.1 & 10.5 & 8.8 & 9.4 & 9.2 \\
Linear reg.+CI & 5.8 & 8.2 & 6.9 & 7.9 & 7.2 \\
Emperor mitigation & 7.8 & 10.2 & 8.2 & 9.4 & 8.9 \\
\bottomrule
\end{tabular}
\end{table}

\begin{table}[t]
\centering
\caption{Interval quality metrics. 4PL+$\gamma$ achieves only 78\% empirical coverage at the nominal 95\% level (undercoverage), with average CI width of 19.5~pp. Linear regression+CI provides both better coverage and narrower intervals.}
\label{tab:interval_quality}
\small
\begin{tabular}{lcccc}
\toprule
Method & Cov.\ 95\% & Cov.\ 80\% & Width 95\% & $\tau$ \\
\midrule
4PL+$\gamma$ & 0.78 & 0.62 & 19.5 & 0.62 \\
Linear reg.+CI & 0.83 & 0.70 & 16.2 & 0.65 \\
Single-sig.\ 2PL+$\gamma$ & 0.80 & 0.65 & 17.8 & 0.60 \\
\bottomrule
\end{tabular}
\end{table}

\begin{table}[t]
\centering
\caption{Dose-response $\gamma$ estimates. The contamination effect is identifiable only at 50\%+ dosage. At 5\% and 25\%, the 95\% CI includes zero, meaning the model cannot distinguish contamination from noise.}
\label{tab:dose_response}
\small
\begin{tabular}{lcccc}
\toprule
Dosage & $\hat{\gamma}$ & 95\% CI & Excl.\ zero? \\
\midrule
5\% & 0.04 & [$-$0.18, 0.26] & No \\
25\% & 0.11 & [$-$0.08, 0.31] & No \\
50\% & 0.22 & [0.01, 0.44] & Yes \\
100\% & 0.38 & [0.12, 0.65] & Yes \\
\bottomrule
\end{tabular}
\end{table}

Tables~\ref{tab:calibration_mae}--\ref{tab:dose_response} present the central negative result: the contamination-augmented 4PL IRT model does not significantly outperform simpler baselines.
Linear regression+bootstrap CI achieves the best average MAE (7.2~pp vs.\ 7.8~pp for 4PL+$\gamma$; paired $t$-test $p = 0.23$), with better interval coverage (83\% vs.\ 78\% at nominal 95\%) and tighter intervals (16.2~pp vs.\ 19.5~pp).

The dose-response analysis (Table~\ref{tab:dose_response}) reveals that $\gamma$ is identifiable only at $\geq$50\% dosage.
The $\gamma$--$\theta$ posterior correlation of 0.52 directly evidences the confounding between ability and contamination: the model cannot reliably distinguish ``the model knows the answer'' from ``the model memorized the answer.''
Stan diagnostics show 3.2\% divergent transitions and $\hat{R}_{\max} = 1.08$, suggesting posterior geometry difficulties consistent with near-unidentifiability.

The dose-response curve is not fully monotonic (Spearman $\rho = 0.55$), with 55\% of conditions yielding $\gamma$ CIs that include zero.
These findings establish that at realistic signal quality levels, IRT model complexity is counterproductive---the additional parameters amplify estimation noise rather than capturing real structure.

% ---- Experiment 5: Noise Propagation ----
\subsection{Noise Propagation Analysis}
\label{sec:noise}

\begin{table}[t]
\centering
\caption{CI width and coverage under noise injection. At matched noise level (1$\times$), 78\% of conditions have CI width exceeding the 15~pp practical threshold. The measurement error model converges but does not resolve the fundamental signal quality problem.}
\label{tab:noise_propagation}
\small
\begin{tabular}{lcccc}
\toprule
Noise & $\gamma$ CI (pp) & $\theta$ CI (pp) & Cov.\ 95\% & Cov.\ 80\% \\
\midrule
Baseline (0$\times$) & 18.0 & 12.0 & 0.82 & 0.68 \\
Matched (1$\times$) & 22.5 & 15.8 & 0.88 & 0.74 \\
1.5$\times$ & 28.0 & 20.5 & 0.93 & 0.82 \\
2$\times$ & 35.2 & 27.1 & 0.97 & 0.90 \\
Meas.\ error model & 24.0 & 17.2 & 0.90 & 0.76 \\
\bottomrule
\end{tabular}
\end{table}

Table~\ref{tab:noise_propagation} quantifies the pathway from signal noise to calibration futility.
Even at baseline (no additional noise), the $\gamma$ CI width is 18.0~pp, already exceeding any reasonable practical threshold.
At matched noise (1$\times$, corresponding to Self-Critique's estimated FP rate of 35\% and FN rate of 28\%), the CI inflates to 22.5~pp, and 78\% of conditions exceed the 15~pp threshold.
Coverage increases with noise (from 0.82 to 0.97) not because calibration improves, but because intervals widen---a model producing a 95\% CI of [$42.5\%, 77.5\%$] for a 60\% accuracy model is technically ``calibrated'' but practically useless.

This analysis yields an operational definition of \emph{practical identifiability}: the regime where CI width $< 15$~pp, which requires signal AUC $> 0.70$ under current noise conditions.

% ---- Experiment 6: Prior Sensitivity ----
\subsection{Prior Sensitivity Analysis}
\label{sec:prior}

\begin{table}[t]
\centering
\caption{Prior sensitivity of $\gamma$ estimation. MAE varies by 2.3~pp (30\% of point estimate) across prior choices, confirming that data are insufficiently informative to override the prior---a hallmark of under-identification.}
\label{tab:prior_sensitivity}
\small
\begin{tabular}{lcccc}
\toprule
Prior & $\hat{\gamma}$ & MAE (pp) & Cov.\ 95\% & Width (pp) \\
\midrule
$\mathcal{N}(0, 0.5)$ & 0.22 & 9.5 & 0.85 & 15.2 \\
$\mathcal{N}(0, 1.0)$ [default] & 0.38 & 7.8 & 0.78 & 19.5 \\
$\mathcal{N}(0, 2.0)$ & 0.51 & 7.2 & 0.72 & 24.8 \\
\bottomrule
\end{tabular}
\end{table}

Table~\ref{tab:prior_sensitivity} demonstrates that calibration results are highly sensitive to the $\gamma$ prior.
The tight prior $\mathcal{N}(0, 0.5)$ shrinks $\hat{\gamma}$ by 45\%, while the wide prior $\mathcal{N}(0, 2.0)$ inflates it by 35\%.
MAE ranges from 7.2 to 9.5~pp across prior choices---a 2.3~pp sensitivity range representing 30\% of the point estimate.
In a well-identified model, the posterior should be dominated by the likelihood rather than the prior; this sensitivity is a direct symptom of weak signal informativeness (prior sensitivity index = 0.38).

% ---- Experiment 7: Few-Examinee Regime ----
\subsection{Few-Examinee Regime}
\label{sec:few_examinee}

\begin{table}[t]
\centering
\caption{Identifiability metrics by number of evaluated models at signal AUC = 0.65. $\gamma$ becomes statistically significant at $N \geq 80$; nominal coverage requires $N \geq 150$. Growing Pains pre-estimation provides only marginal improvement ($\leq$15\% at $N=20$) because the bottleneck is $\gamma$--$\theta$ confounding, not item parameter estimation.}
\label{tab:few_examinee}
\small
\begin{tabular}{lccccc}
\toprule
$N$ & $\gamma$ RMSE & Rel.\ bias & $\theta$ corr.\ & Cov.\ 95\% & GP $\Delta$ \\
\midrule
20 & 0.42 & 55\% & 0.58 & 68\% & +15\% \\
50 & 0.28 & 38\% & 0.72 & 75\% & +10\% \\
100 & 0.18 & 25\% & 0.82 & 80\% & +5\% \\
200 & 0.12 & 18\% & 0.88 & 85\% & +3\% \\
\bottomrule
\end{tabular}
\end{table}

Table~\ref{tab:few_examinee} presents a simulation study of identifiability as a function of model count, fixing signal AUC at 0.65 (the level observed in our controlled experiments).
At $N = 20$---a common scenario for benchmark evaluation---$\gamma$ relative bias reaches 55\% and coverage drops to 68\%.
Even at $N = 200$, relative bias remains at 18\% and coverage at 85\%, still below the nominal 95\%.
The Growing Pains pre-estimation strategy~\citep{bhattpain}, which borrows item parameter priors from a larger pool, provides diminishing returns ($+3\%$ at $N = 200$) because the bottleneck is $\gamma$--$\theta$ joint estimation, which cannot be resolved by better item priors alone.

The minimum model count for $\gamma$ statistical significance is $N = 80$; for nominal coverage, $N = 150$.
These thresholds, combined with the signal AUC requirement of $\geq 0.70$, define the joint identifiability boundary.

% ---- Experiment 8: Signal Survival ----
\subsection{Signal Survival Analysis}
\label{sec:survival}

\begin{table}[t]
\centering
\caption{Detection signal AUC across training regimes. RL degrades all signals below the 0.70 identifiability threshold (average survival rate 73\%). The best RL-regime signal (Self-Critique, 0.57) provides negligible discrimination. ANOVA confirms a significant regime effect ($p = 0.002$).}
\label{tab:signal_survival}
\small
\begin{tabular}{lcccc}
\toprule
Signal & Base & SFT-Clean & SFT-Contam & RL \\
\midrule
Embedding Retrieval & 0.78 & 0.72 & 0.72 & 0.60 \\
Min-K\%++ & 0.73 & 0.65 & 0.68 & 0.55 \\
ConStat & 0.71 & 0.63 & 0.66 & 0.53 \\
Self-Critique & 0.60 & 0.58 & 0.63 & 0.57 \\
\bottomrule
\end{tabular}
\end{table}

\begin{table}[t]
\centering
\caption{Signal survival rate (AUC retention relative to base model). Self-Critique shows paradoxical stability (95\% retention under RL) only because it starts from a low baseline (0.60). Post-training degradation is systematic, not signal-specific.}
\label{tab:survival_rate}
\small
\begin{tabular}{lccc}
\toprule
Signal & SFT-Clean & SFT-Contam & RL \\
\midrule
Embedding Retrieval & 0.92 & 0.92 & 0.77 \\
Min-K\%++ & 0.89 & 0.93 & 0.75 \\
ConStat & 0.89 & 0.93 & 0.75 \\
Self-Critique & 0.97 & 1.05 & 0.95 \\
\bottomrule
\end{tabular}
\end{table}

Tables~\ref{tab:signal_survival}--\ref{tab:survival_rate} present a core finding: detection signals degrade systematically with post-training.
SFT reduces signal AUC by an average of 12\%, while RL causes a 27\% reduction.
Under RL, \emph{all} signals fall below the 0.70 identifiability threshold established in our simulation study, with the best signal (Self-Critique at 0.57) providing marginal discrimination at best.

Two critical observations emerge.
First, the AUC difference between SFT-clean and SFT-contaminated regimes is only 0.03 on average, indicating that signal degradation is driven by the post-training \emph{process itself}, not by contamination---signals detect ``training change'' rather than ``contamination change.''
Second, Self-Critique's apparent stability under RL (95\% retention) is a pyrrhic victory: it retains AUC in the 0.57--0.63 range, which is consistently below the identifiability threshold.

% ---- Experiment 9: Scale Bridging ----
\subsection{Scale Analysis: 7B and 72B}
\label{sec:scale}

\begin{table}[t]
\centering
\caption{Calibration degrades with model scale. Comparing controlled results at 100\% dosage, 7B shows weaker signal AUC (0.66 vs.\ 0.69), higher $\gamma$--$\theta$ posterior correlation (0.58 vs.\ 0.52), and worse calibration MAE (8.5~pp vs.\ 7.8~pp) relative to 3B.}
\label{tab:7b_comparison}
\small
\begin{tabular}{lccc}
\toprule
Metric & Qwen-3B & Qwen-7B \\
\midrule
MMLU inflation (pp) & 5.2 & 4.5 \\
Best signal AUC & 0.69 & 0.66 \\
$\hat{\gamma}$ & 0.34 & 0.25 \\
$\gamma$--$\theta$ posterior corr.\ & 0.52 & 0.58 \\
Calibration MAE (pp) & 7.8 & 8.5 \\
CI width 95\% (pp) & 19.5 & 22.0 \\
\bottomrule
\end{tabular}
\end{table}

\begin{table}[t]
\centering
\caption{Observational analysis at 72B scale. Signal AUC (best: 0.60) and calibration MAE (11.2~pp) confirm the inverse relationship between model scale and calibration feasibility, extrapolating the trend from controlled experiments.}
\label{tab:72b_comparison}
\small
\begin{tabular}{lccc}
\toprule
Metric & Qwen-3B & Qwen-7B & Qwen-72B \\
\midrule
Best signal AUC & 0.69 & 0.66 & 0.60 \\
RRF AUC & 0.67 & 0.63 & 0.59 \\
$\hat{\gamma}$ & 0.34 & 0.25 & 0.04 \\
$\gamma$ CI excl.\ zero & Yes & Marginal & No \\
Calibration MAE (pp) & 7.8 & 8.5 & 11.2 \\
\bottomrule
\end{tabular}
\end{table}

Tables~\ref{tab:7b_comparison}--\ref{tab:72b_comparison} demonstrate that identifiability degrades monotonically with model scale.
At 7B, inflation per unit dosage decreases (4.5~pp vs.\ 5.2~pp at 3B) because stronger pretraining knowledge reduces the marginal benefit of memorization.
This reduced contamination signal-to-noise ratio manifests as lower detection AUC (0.66), higher $\gamma$--$\theta$ confounding (0.58), and worse calibration MAE (8.5~pp).

At 72B (observational, without controlled injection), the pattern reaches its extreme: the best signal AUC is 0.60, $\hat{\gamma} = 0.04$ with a CI that includes zero, and calibration MAE is 11.2~pp with CI width of 26.0~pp.
This creates a \emph{fundamental scale barrier}: the models most scrutinized for contamination (large, frontier models) are precisely those where calibration is least feasible, because ``contamination'' and ``knowledge'' become conceptually indistinguishable at sufficient scale.

% ---- Experiment 10: RL Observational Case Study ----
\subsection{RL Observational Case Study}
\label{sec:rl_case}

\begin{table}[t]
\centering
\caption{Calibration on RL-trained models (8 models including Qwen2.5-Instruct and DeepSeek-R1-Distill). 4PL+$\gamma$ \emph{increases} MAE relative to uncalibrated scores (10.5 vs.\ 9.8~pp), with negative expected regret reduction ($-$0.03), demonstrating that calibration can actively harm model selection when signal quality is insufficient.}
\label{tab:rl_case_study}
\small
\begin{tabular}{lcccc}
\toprule
Method & MAE & CI Width & Cov.\ 95\% & Regret $\Delta$ \\
\midrule
Uncalibrated & 9.8 & --- & --- & --- \\
4PL+$\gamma$ & 10.5 & 24.5 & 0.82 & $-$0.03 \\
Remove-rescore & 9.2 & --- & --- & --- \\
Linear reg.+CI & 9.9 & 21.0 & 0.85 & +0.01 \\
\bottomrule
\end{tabular}
\end{table}

\begin{table}[t]
\centering
\caption{Per-model $\gamma$ estimates for RL-trained models. 85\% of conditions yield CIs that include zero. Contamination is statistically undetectable in the RL regime. Specificity check: $\hat{\gamma} = 0.02$ on a clean benchmark confirms no systematic false positives.}
\label{tab:rl_gamma}
\small
\begin{tabular}{lcccc}
\toprule
Model & $\hat{\gamma}$ & 95\% CI lower & 95\% CI upper & Excl.\ zero \\
\midrule
Qwen2.5-0.5B-Inst & 0.12 & $-$0.15 & 0.39 & No \\
Qwen2.5-1.5B-Inst & 0.08 & $-$0.18 & 0.34 & No \\
Qwen2.5-3B-Inst & 0.05 & $-$0.21 & 0.31 & No \\
Qwen2.5-7B-Inst & 0.03 & $-$0.25 & 0.31 & No \\
DS-R1-Distill-1.5B & 0.10 & $-$0.14 & 0.34 & No \\
DS-R1-Distill-7B & 0.06 & $-$0.20 & 0.32 & No \\
\bottomrule
\end{tabular}
\end{table}

Tables~\ref{tab:rl_case_study}--\ref{tab:rl_gamma} present the most striking result: in the RL regime, calibration is not merely unhelpful---it is \emph{harmful}.
The 4PL+$\gamma$ model produces higher MAE than uncalibrated scores (10.5 vs.\ 9.8~pp) with negative regret reduction ($-$0.03), meaning calibration degrades model selection decisions.
This occurs because noisy exposure scores (RRF AUC~0.59) cause the IRT model to ``hallucinate'' contamination effects, overcorrecting false positives while undercorrecting true positives.

For 85\% of RL model-benchmark conditions, $\gamma$ CIs include zero.
This is not because RL models are uncontaminated---it is because the signals are too noisy to detect any contamination that may exist.
Top-5 model selection agreement with ground truth is only 60\%, barely above chance for a pool of 8 models.
The specificity check on a clean benchmark ($\hat{\gamma} = 0.02$) confirms that the model does not systematically produce false contamination detections---a low bar that does not redeem the fundamental inability to detect real contamination.

% ---- Experiment 11: Closed-Weight 2-Signal Mode ----
\subsection{Closed-Weight 2-Signal Degraded Mode}
\label{sec:closed_weight}

\begin{table}[t]
\centering
\caption{Calibration quality gap between 4-signal (open-weight) and 2-signal (closed-weight) modes. The 2-signal mode suffers 38\% higher MAE and Kendall $\tau$ drops below 0.50, meaning calibrated rankings are worse than random for model comparison.}
\label{tab:closed_weight}
\small
\begin{tabular}{lccc}
\toprule
Metric & 4-signal & 2-signal & Degradation \\
\midrule
RRF AUC (100\%) & 0.70 & 0.64 & $-$8.6\% \\
MAE (pp) & 7.8 & 10.8 & +38\% \\
CI Width 95\% (pp) & 19.5 & 25.2 & +29\% \\
Kendall $\tau$ & 0.62 & 0.48 & $-$23\% \\
Coverage 95\% & 0.78 & 0.75 & $-$3.8\% \\
\bottomrule
\end{tabular}
\end{table}

Table~\ref{tab:closed_weight} quantifies the gap between open-weight (4-signal) and closed-weight (2-signal: ConStat + Self-Critique) calibration.
The 2-signal mode shows 38\% higher MAE (10.8 vs.\ 7.8~pp) and Kendall $\tau$ drops to 0.48---below 0.50, meaning calibrated model rankings are not consistently better than uncalibrated ones.
The $\gamma$ CI includes zero for 72\% of conditions.

This result carries significant practical implications: closed-weight models (GPT-4, Claude, Gemini)---which are the most scrutinized for contamination and the primary targets for calibration---are precisely the models for which calibration is least feasible.
The available signals (ConStat and Self-Critique) are the two weakest of the four, and their fusion AUC of 0.64 falls well below the 0.70 identifiability threshold.

% ============================================================
\section{Discussion: The Identifiability Map}
\label{sec:discussion}

\begin{table}[t]
\centering
\caption{Synthesized identifiability boundaries. Each row is a condition that must be met for contamination-aware IRT calibration to be feasible. The ``Typical'' column indicates whether the condition is commonly satisfied in current LLM evaluation practice. Only 1 of 5 conditions is typically met.}
\label{tab:identifiability_map}
\small
\begin{tabular}{lcc}
\toprule
Condition & Threshold & Typical? \\
\midrule
Signal AUC & $\geq 0.70$ & Rarely \\
Model count $N$ & $\geq 80$ & Sometimes \\
CI width & $< 15$~pp & Rarely \\
Post-training regime & SFT only & Decreasingly \\
Model scale & $\leq 3$B & No \\
\bottomrule
\end{tabular}
\end{table}

Our experiments collectively delineate an \textbf{identifiability map} with three dimensions: signal quality, model count, and model scale (Table~\ref{tab:identifiability_map}).
The viable region---where calibration produces actionable results---is defined by the joint requirement: signal AUC $\geq 0.70$, $N \geq 80$, CI width $< 15$~pp, SFT-only post-training, and model scale $\leq 3$B.
In current LLM evaluation practice, these conditions are rarely simultaneously satisfied.

\paragraph{The ability--contamination confound.}
The $\gamma$--$\theta$ posterior correlation increases with model scale (0.52 at 3B, 0.58 at 7B), reflecting a fundamental conceptual challenge: for large models trained on internet-scale data, the distinction between ``knowledge acquired from contamination'' and ``knowledge acquired from legitimate training'' becomes blurred.
This is not a limitation of our method but of the calibration \emph{problem formulation} itself.

\paragraph{Signal quality as the binding constraint.}
Across all experiments, the single most important factor is detection signal AUC.
When AUC $< 0.70$, noise propagation renders all downstream calibration methods---regardless of complexity---practically useless.
The catch-22 identified in Section~\ref{sec:signal_quality} (signals strongest where calibration is least needed) suggests that improving signal quality at low contamination levels is the most impactful direction for future work.

\paragraph{When to calibrate.}
Our results yield a concrete decision rule: \emph{attempt calibration only when} the best available detection signal achieves AUC $\geq 0.70$ on a validation set, $N \geq 80$ models are available, the models are SFT-trained (not RL), and model scale is $\leq 7$B.
When these conditions are not met, the honest answer is to report uncalibrated scores with an explicit contamination warning rather than to apply calibration that may be harmful (Section~\ref{sec:rl_case}).

\paragraph{Implications for detection research.}
The identifiability map provides concrete targets for detection signal development: AUC $\geq 0.80$ would expand the feasible region substantially, particularly if achieved at low contamination dosage ($\leq 25\%$) and for RL-trained models.
Furthermore, our aggregation analysis (Section~\ref{sec:rrf}) suggests that new signals should aim for low correlation ($r < 0.30$) with existing signals to enable effective fusion.

% ============================================================
\section{Conclusion}

We set out to build a contamination-aware calibration system and instead discovered the conditions under which such a system can work.
Our systematic investigation---spanning controlled SFT injection across four model scales, four benchmarks, and four detection signals---establishes empirical identifiability boundaries for contamination-aware IRT calibration.
The core finding is sobering: calibration is feasible only under conditions rarely met in practice (signal AUC $\geq 0.70$, $N \geq 80$, small-scale SFT models), and actively harmful when applied outside these boundaries.

Rather than viewing these as purely negative results, we frame them as \emph{necessary groundwork}: before the field can build reliable calibration systems, it must understand when calibration is solvable.
Our identifiability map provides concrete, quantitative targets for future detection research and establishes a principled criterion for when to attempt calibration versus when to honestly report the limits of current methodology.

% ============================================================
\bibliographystyle{aaai}

\begin{thebibliography}{10}

\bibitem{yang2023rethinking}
S.~Yang, W.~Chiang, L.~Zheng, J.~Gonzalez, and I.~Stoica.
\newblock Rethinking benchmark and contamination for language models with rephrased samples.
\newblock {\em arXiv:2311.04850}, 2023.

\bibitem{oren2025emperor}
Y.~Oren, N.~Meister, N.~Chatterji, F.~Ladhak, and T.~Hashimoto.
\newblock Proving test set contamination in black-box language models.
\newblock In {\em ICML}, 2025.

\bibitem{eckerly2017detecting}
C.~Eckerly.
\newblock Detecting item preknowledge and item compromise.
\newblock {\em Journal of Educational Measurement}, 54(2):124--142, 2017.

\bibitem{sinharay2017detection}
S.~Sinharay.
\newblock Detection of item preknowledge using likelihood ratio test and score test.
\newblock {\em Journal of Educational and Behavioral Statistics}, 42(1):46--68, 2017.

\bibitem{dodge2021documenting}
J.~Dodge, M.~Sap, A.~Marasovic, W.~Agnew, G.~Ilharco, D.~Groeneveld, M.~Mitchell, and M.~Gardner.
\newblock Documenting large webtext corpora: A case study on the Colossal Clean Crawled Corpus.
\newblock In {\em EMNLP}, 2021.

\bibitem{shi2024detecting}
W.~Shi, A.~Ajith, M.~Xia, Y.~Huang, D.~Liu, T.~Blevins, D.~Chen, and L.~Zettlemoyer.
\newblock Detecting pretraining data from large language models.
\newblock In {\em ICLR}, 2024.

\bibitem{dekoninck2024constat}
J.~Dekoninck, M.~Müller, M.~Fischer, and L.~Vechev.
\newblock ConStat: Performance-based contamination detection in large language models.
\newblock In {\em NeurIPS}, 2024.

\bibitem{weller2024selfcritique}
O.~Weller, M.~Marone, V.~Braverman, D.~Lawrie, and B.~Van~Durme.
\newblock According to \ldots: Prompting language models improves quoting from pre-training data.
\newblock In {\em EACL}, 2024.

\bibitem{polo2024tinybenchmarks}
F.~Polo, L.~Weber, L.~Choshen, Y.~Sun, G.~Xu, and M.~Yurochkin.
\newblock tinyBenchmarks: Evaluating LLMs with fewer examples.
\newblock In {\em ICML}, 2024.

\bibitem{bhattpain}
S.~Bhatt, A.~Ladhak, and B.~Van~Durme.
\newblock Growing pains: Understanding the limitations of scaling IRT for LLM evaluation.
\newblock {\em arXiv preprint}, 2024.

\bibitem{zhuang2024atlas}
Y.~Zhuang, Y.~Liu, Q.~Ning, and W.~Wang.
\newblock ATLAS: Adaptive item selection for LLM evaluation.
\newblock {\em arXiv preprint}, 2024.

\bibitem{paganini2024hibayes}
M.~Paganini.
\newblock HiBayES: Hierarchical Bayesian evaluation of LLMs.
\newblock {\em arXiv preprint}, 2024.

\end{thebibliography}

\end{document}
